\subsection{Xây Dựng Mô Hình (Modeling)}


Bốn mô hình tuyến tính chính được so sánh:
\begin{enumerate}
    \item \textbf{OLS full}: dùng toàn bộ 11 feature sau tiền xử lý.
    \item \textbf{OLS selected}: chọn feature có p-value nhỏ hơn 0.05 từ OLS full.
    \item \textbf{Ridge}: dùng hàm \texttt{ridge\_fit} ở Phần 1 và chọn $\lambda$ bằng 5-fold CV.
    \item \textbf{Lasso}: dùng hàm \texttt{lasso\_fit} ở Phần 1 và chọn $\lambda$ bằng 5-fold CV.
\end{enumerate}

\subsubsection{OLS full và chọn biến theo p-value}

OLS full có $R^2_{\mathrm{train}} = 0.681784$ và
$R^2_{\mathrm{adj,train}} = 0.677690$. Các biến có ý nghĩa thống kê ở mức 0.05 được giữ lại cho
OLS selected:
\begin{center}
\texttt{pl\_trandur}, \texttt{pl\_imppar}, \texttt{st\_teff},
\texttt{log\_pl\_trandep}, \texttt{log\_pl\_bmasse}, \texttt{log\_sy\_dist}.
\end{center}

\begin{table}[H]
\centering
\caption{Bảng suy diễn hệ số OLS full}
\label{tab:part2_ols_inference}
\resizebox{\textwidth}{!}{%
\begin{tabular}{@{}lrrrrrr@{}}
\hline
\textbf{Feature} & \textbf{Coef} & \textbf{Std.Err.} & \textbf{t-stat} & \textbf{p-value} &
\textbf{CI lower} & \textbf{CI upper} \\
\hline
\texttt{pl\_orbeccen} & -0.004508 & 0.002704 & -1.6673 & 0.095816 & -0.009815 & 0.000799 \\
\texttt{pl\_trandur} & -0.011414 & 0.004214 & -2.7088 & 0.006888 & -0.019684 & -0.003144 \\
\texttt{pl\_imppar} & 0.005751 & 0.002482 & 2.3169 & 0.020743 & 0.000879 & 0.010622 \\
\texttt{st\_teff} & -0.014611 & 0.007290 & -2.0043 & 0.045348 & -0.028918 & -0.000303 \\
\texttt{st\_rad} & 0.009279 & 0.005864 & 1.5822 & 0.113964 & -0.002231 & 0.020790 \\
\texttt{st\_met} & 0.000139 & 0.002609 & 0.0533 & 0.957505 & -0.004982 & 0.005260 \\
\texttt{log\_pl\_orbsmax} & 0.007766 & 0.005142 & 1.5104 & 0.131315 & -0.002326 & 0.017857 \\
\texttt{log\_pl\_trandep} & 0.013692 & 0.003279 & 4.1757 & $3.28\times10^{-5}$ & 0.007256 & 0.020128 \\
\texttt{log\_pl\_insol} & -0.005905 & 0.005809 & -1.0165 & 0.309668 & -0.017305 & 0.005496 \\
\texttt{log\_pl\_bmasse} & 0.095317 & 0.002466 & 38.6447 & $8.78\times10^{-190}$ & 0.090476 & 0.100158 \\
\texttt{log\_sy\_dist} & 0.034443 & 0.004020 & 8.5689 & $4.84\times10^{-17}$ & 0.026554 & 0.042332 \\
\hline
\end{tabular}}
\end{table}

Biến \texttt{log\_pl\_bmasse} có hệ số dương lớn nhất và p-value cực nhỏ, phù hợp với trực giác
vật lý: hành tinh có khối lượng lớn thường có bán kính lớn hơn.

\subsubsection{Chọn siêu tham số cho Ridge và Lasso}

Ridge quét lưới $\lambda$ từ $10^{-3}$ đến $10^3$ theo thang log. Giá trị tốt nhất theo CV-MSE là:
\begin{equation}
    \lambda^*_{\mathrm{Ridge}} = 10.
\end{equation}

\begin{table}[H]
\centering
\caption{Một số kết quả Ridge CV}
\label{tab:part2_ridge_cv}
\begin{tabular}{@{}rrrr@{}}
\hline
\textbf{$\lambda$} & \textbf{Mean CV-MSE} & \textbf{Std CV-MSE} & \textbf{Mean CV-$R^2$} \\
\hline
0.0010 & 0.005075 & 0.001660 & 0.673407 \\
0.1000 & 0.005075 & 0.001659 & 0.673410 \\
1.0000 & 0.005075 & 0.001656 & 0.673433 \\
3.1623 & 0.005074 & 0.001650 & 0.673469 \\
10.0000 & 0.005074 & 0.001629 & 0.673433 \\
31.6228 & 0.005090 & 0.001571 & 0.672223 \\
100.0000 & 0.005240 & 0.001432 & 0.661976 \\
1000.0000 & 0.008594 & 0.001069 & 0.441426 \\
\hline
\end{tabular}
\end{table}

Sau khi mở rộng lưới, Lasso quét 13 giá trị
$\lambda \in [10^{-4}, 10^2]$ theo thang log bằng 5-fold CV. Lưới mới cho thấy vùng
regularization tốt khá phẳng ở các $\lambda$ nhỏ, sau đó CV-MSE bắt đầu tăng rõ khi
$\lambda \ge 0.3162$. Giá trị tốt nhất là:
\begin{equation}
    \lambda^*_{\mathrm{Lasso}} = 0.1.
\end{equation}

\begin{table}[H]
\centering
\caption{Kết quả Lasso CV}
\label{tab:part2_lasso_cv}
\small
\renewcommand{\arraystretch}{1.0}
\begin{adjustbox}{max width=\textwidth}
\begin{tabular}{@{}rrrr@{}}
\hline
\textbf{$\lambda$} & \textbf{Mean CV-MSE} & \textbf{Std CV-MSE} & \textbf{Mean CV-$R^2$} \\
\hline
0.0001 & 0.005075 & 0.001659 & 0.673409 \\
0.0003 & 0.005075 & 0.001659 & 0.673409 \\
0.001 & 0.005075 & 0.001659 & 0.673410 \\
0.003 & 0.005075 & 0.001659 & 0.673413 \\
0.010 & 0.005075 & 0.001659 & 0.673421 \\
0.032 & 0.005075 & 0.001658 & 0.673442 \\
0.100 & 0.005074 & 0.001655 & 0.673462 \\
0.316 & 0.005081 & 0.001648 & 0.672981 \\
1.000 & 0.005105 & 0.001628 & 0.671292 \\
3.162 & 0.005249 & 0.001575 & 0.661616 \\
10.000 & 0.005782 & 0.001438 & 0.626800 \\
31.623 & 0.007753 & 0.001154 & 0.497348 \\
100.000 & 0.015461 & 0.001291 & -0.007357 \\
\hline
\end{tabular}
\end{adjustbox}
\end{table}

Với $\lambda^* = 0.1$, Lasso vẫn không triệt tiêu hệ số nào: $0/11$ feature có hệ số bằng
0 theo ngưỡng $10^{-10}$. Điều này cho thấy Lasso trong lần chạy này hoạt động giống một OLS được
làm mượt nhẹ hơn là một bộ chọn biến thưa. Lưới mở rộng cũng làm rõ rằng nếu phạt mạnh hơn
($\lambda \ge 1$), sai số CV tăng dần; do đó việc không chọn nghiệm thưa là kết quả thực nghiệm
của dữ liệu, không phải do lưới ban đầu quá hẹp.


